
\begin{center}
    {\LARGE \textbf{Базовые понятия комбинаторики}}\\
    \vspace{4pt}
    {\small Правило сложения/умножения, перестановки, размещения без повторений, сочетания без повторений}
\end{center}

\section{Правило сложения}

Если некоторый результат можно получить \textbf{либо} способом типа $A$ (число вариантов $m$),
\textbf{либо} способом типа $B$ (число вариантов $n$), и эти способы \textbf{взаимно исключают} друг друга,
то общее число способов равно
\[
m+n.
\]

\textbf{Пример.} На полке $5$ книг по алгебре и $3$ по геометрии. Выбрать одну книгу можно
\[
5+3=8
\]
способами.

\newpage

\hspace{3mm}

\section{Правило умножения}

Если действие выполняется в несколько \textbf{последовательных} шагов и:
\begin{itemize}
    \item первый шаг можно выполнить $m$ способами;
    \item второй шаг можно выполнить $n$ способами (для каждого варианта первого шага);
\end{itemize}
то общее число способов выполнить оба шага равно
\[
m\cdot n.
\]

\textbf{Пример.} Есть $3$ варианта рубашек и $4$ варианта брюк. Число комплектов:
\[
3\cdot 4 = 12.
\]

\newpage

\hspace{3mm}

\section{Перестановки}

\textbf{Перестановка} из $n$ различных элементов — это любое упорядочивание \textbf{всех} $n$ элементов.

Обозначение: $P_n$.

Формула:
\[
P_n = n!,
\qquad
n! = 1\cdot 2\cdot 3\cdots n.
\]

\textbf{Пример.} Рассадить $4$ человек по местам можно
\[
4! = 24
\]
способами.


\newpage

\hspace{3mm}

\section{Размещения без повторений (порядок важен)}

\textbf{Размещение} из $n$ по $k$ (без повторений) — это выбор $k$ элементов из $n$,
где \textbf{важен порядок}, и каждый элемент можно использовать \textbf{не более одного раза}.

Обозначение: $A_n^k$.

Формулы (эквивалентные записи):
\[
A_n^k = n(n-1)(n-2)\cdots(n-k+1)
\]
и
\[
A_n^k = \frac{n!}{(n-k)!}.
\]

\textbf{Пример.} Выбрать и упорядочить $2$ книги из $5$ можно
\[
A_5^2=\frac{5!}{3!}=\frac{120}{6}=20
\]
способами.


\newpage

\hspace{3mm}

\section{Сочетания без повторений (порядок не важен)}

\textbf{Сочетание} из $n$ по $k$ (без повторений) — это выбор $k$ элементов из $n$,
где \textbf{порядок не важен}, и каждый элемент используется \textbf{не более одного раза}.

Обозначение: $C_n^k$ (в англоязычной литературе пишут $\binom{n}{k}$, но мы так писать не будем!). Читается как <<цэ из эн по ка>>.

Формула:
\[
C_n^k = \binom{n}{k} = \frac{n!}{k!(n-k)!}.
\]

Связь с размещениями:
\[
C_n^k = \frac{A_n^k}{k!}.
\]

\textbf{Пример.} Выбрать $2$ книги из $5$ можно
\[
C_5^2=\frac{5!}{2!\,3!}=\frac{120}{12}=10
\]
способами.


\newpage

\hspace{3mm}

\section{Сравнение формул}

\begin{table}[h]
\begin{tabular}{|l|l|l|l|}
\hline Тип конструкции & Порядок важен? & Повторения & Формула \\
\hline Перестановки & Да & Нет & $n!$ \\
\hline Размещения $A_n^k$ & Да & Нет & $\frac{n!}{(n-k)!}$ \\
\hline Сочетания $C_n^k$ & Нет & Нет & $\frac{n!}{k!(n-k)!}$ \\
\hline
\end{tabular}
\end{table}


\newpage

\hspace{3mm}

\section{Как быстро выбрать нужную формулу}

\begin{itemize}
    \item Используются \textbf{все} элементы и важен порядок $\Rightarrow$ \textbf{перестановки} $P_n$.
    \item Берём \textbf{часть} элементов и важен порядок $\Rightarrow$ \textbf{размещения} $A_n^k$.
    \item Берём \textbf{часть} элементов и порядок \textbf{не важен} $\Rightarrow$ \textbf{сочетания} $C_n^k$.
\end{itemize}
